\chapter{Stop-loss route}\label{chap:stoploss}

This chapter develops the one-dimensional mechanism behind the comparison.
The convex-order question is translated into the stop-loss transform
$u \mapsto \E[(X-u)_+]$: a sharp universal envelope for that transform is
proved under the sub-Gaussian tail cap, and the envelope bound is then
lifted back to full convex domination.

\section{Layer-cake and stop-loss identities}

\begin{lemma}[Layer-cake for hinge functions]\label{lem:layercake}
\lean{layer_cake_hinge}\leanok
For every probability measure $\mu$ on $\R$ and every $u \ge 0$,
\[
  \int (x-u)_+\,d\mu(x) = \int_u^\infty \Prb_\mu(x > t)\,dt.
\]
\end{lemma}

\begin{proof}\leanok
Expand $(x-u)_+$ as $\int_u^\infty \mathbf{1}_{\{x>t\}}\,dt$ and apply
Tonelli's theorem to swap the order of integration.
\end{proof}

\begin{lemma}[Stop-loss as a tail integral]\label{lem:stoploss-tail}
\lean{stopLoss_eq_integral_tail}\leanok
If $\mu$ is a probability measure with finite first moment, then for every
$u \ge 0$ its stop-loss transform equals the tail integral from
Lemma~\ref{lem:layercake}.
\end{lemma}

\begin{proof}\leanok
Immediate from Lemma~\ref{lem:layercake} and the definition of the
stop-loss transform.
\end{proof}

\section{The sharp stop-loss envelope}

\begin{theorem}[Sharp stop-loss envelope under the sub-Gaussian cap]\label{thm:envelope}
\lean{stopLoss_envelope}\leanok
Assume that $\mu$ is centered and satisfies the tail cap
$\Prb_\mu(|x|>t)\le s_G(t)$ for every $t \ge 0$. Then for every $u \ge 0$,
\[
  \mathrm{stopLoss}_\mu(u) \le J_G(u),
\]
where $J_G$ is the explicit envelope built from the constants
$B$, $a$, and $p_0$.
\end{theorem}

\begin{proof}\leanok
Split into the two regions $u \le a$ and $u \ge a$.
For $u \le a$, the mean-zero condition and the tail cap bound the positive part
of the first moment by $B$, giving the affine upper bound
$B - p_0 u$. For $u \ge a$, the layer-cake identity rewrites the stop-loss as a
tail integral and the cap is integrated directly.
\end{proof}

\begin{theorem}[Extension to all thresholds]\label{thm:extend-all-u}
\lean{stopLoss_extend_all_u}\leanok
The stop-loss bound extends from $u \ge 0$ to every real threshold $u$ by
applying the same argument to $-X$ and using the mean-zero identity relating
the stop-loss transforms of $X$ and $-X$.
\end{theorem}

\begin{proof}\leanok
The negative-threshold case is rewritten in terms of the
positive-threshold stop-loss of $-X$. Because the tail assumption is symmetric
and the mean is zero, the transformed problem satisfies the same hypotheses.
\end{proof}

\section{From stop-loss control to convex domination}

\begin{theorem}[Convex domination from stop-loss comparison]\label{thm:cx-route}
\lean{convexDomination_c0_via_reduction}\leanok
\uses{thm:extend-all-u}
If the Gaussian stop-loss transform dominates the envelope $J_G$ pointwise,
then every centered sub-Gaussian law is convex-dominated by the Gaussian
comparator at that same scale.
\end{theorem}

\begin{proof}\leanok
First prove the comparison for simple convex functions (finite sums of
affine functions and hinges). Then approximate an arbitrary convex integrable
function from below by a monotone sequence of such simple convex functions.
Monotone convergence passes the inequality to the limit.
\end{proof}
